% Generated draft: requires mathematical and source review.
\documentclass[11pt]{article}
\usepackage[margin=25mm]{geometry}
\usepackage{fontspec}
\setmainfont{Latin Modern Roman}
\usepackage{amsmath,amssymb,mathtools,unicode-math}
\setmathfont{Latin Modern Math}
\usepackage{enumitem,needspace,xurl,hyperref}
\setlength{\parindent}{0pt}
\setlength{\parskip}{5pt}
\setlength{\emergencystretch}{3em}
\newcommand{\N}{\mathbb{N}}
\newcommand{\Z}{\mathbb{Z}}
\newcommand{\Q}{\mathbb{Q}}
\newcommand{\R}{\mathbb{R}}
\newcommand{\C}{\mathbb{C}}
\newcommand{\F}{\mathbb{F}}
\newcommand{\A}{\mathbb{A}}
\newcommand{\PP}{\mathbb{P}}
\newcommand{\E}{\mathbb{E}}
\newcommand{\eps}{\varepsilon}
\newcommand{\dd}{\,\mathrm{d}}
\begin{document}
\section*{UnsolvedMath ID: 800776}
\textbf{Source problem code:} catalogue entry
\subsection{Definitions and mathematical statement}
Write \(\phi\colon X\dashrightarrow Y\) for the rational map referred to in the supplied statement. A holomorphic morphism \(f\colon X\to Y'\) is required to be defined everywhere on \(X\). A birational map between varieties is a rational map admitting a rational inverse.

The supplied question is: under the previously stated assumption on the relative anti-canonical divisor, does a structure theorem hold? More specifically, does there exist a holomorphic morphism
\[
f\colon X\longrightarrow Y'
\]
that is birationally equivalent to \(\phi\colon X\dashrightarrow Y\)?

The referenced assumption on the relative anti-canonical divisor is absent from the supplied record. The record also does not specify the hypotheses on \(X\) and \(Y\), the properties of \(\phi\), the permitted targets \(Y'\), or the intended meaning of birational equivalence of these maps. In particular, it does not state whether birational equivalence must be realized solely by changing the target while keeping \(X\) fixed. Consequently, this is an incomplete formulation of the source question: its hypotheses and equivalence requirement must be recovered before a precise quantified mathematical statement can be given.
\subsection{Short English statement}
The question asks whether an unspecified assumption on a relative anti-canonical divisor yields a structure theorem, specifically an everywhere-defined holomorphic map from the original \(X\) that is birationally equivalent to the given rational map. The necessary preceding assumptions and the meaning of this equivalence are missing.
\subsection{Sources}
\begin{itemize}
\item Local UnsolvedMath registry: permanent ID \(800776\), source code \texttt{catalogue entry}, title ``Positively Curved Projective Varieties --- Problem 11.'' The supplied record attributes the problem to Matsumura, with proposed year 2022, and labels it open; this status is supplied metadata, not an independently verified mathematical assessment.
\item Supplied source list: Matsumura, \emph{Open Problems on Structure of Positively Curved Projective Varieties} (2022), Problem 11, \url{https://arxiv.org/abs/2211.09153}. The supplied provenance records access on 2026-10-16 and extraction by \texttt{arxiv\_latex}. These are supplied provenance details; the source has not been independently consulted for this draft.
\end{itemize}
\end{document}
